\chapter{Software Process Model and Project Planning}

\section{Choice of Process Model}

The project followed an \textbf{iterative, sprint-based (Agile) model} rather than a
single waterfall pass. Four properties of the problem made the choice almost
forced.

\begin{enumerate}
  \item \textbf{The hardware was the least predictable part.} Servo power, BLE
        behaviour on iOS and Wi-Fi reliability could only be characterised by
        building and running the gate. A plan that fixed every detail up front
        would have been rewritten after the first bench test.
  \item \textbf{Correctness was cheap to check but expensive to guess.} The session
        lifecycle has a small number of states and many illegal transitions. Writing
        the transitions down as tests and then iterating until they passed was faster
        and safer than elaborating them on paper.
  \item \textbf{The work could be sliced vertically.} Every increment could deliver a
        working path through all three surfaces --- for example, ``start a session
        and see it on the dashboard'' --- instead of completing one layer at a time.
        Vertical slices kept the system demonstrable after every sprint.
  \item \textbf{The team was small and part-time.} Short sprints with a single
        demonstrable outcome tolerate academic schedules far better than a monolithic
        phase plan.
\end{enumerate}

The model was Scrum-inspired but deliberately lightweight: a single prioritised
backlog derived from the architecture specification, one-week to ten-day sprints, a
working demonstration at the end of each sprint, and a written decision log when a
question could not be settled by the specification. There were no daily stand-ups
and no story-point estimates; progress was measured by what could be executed.

The definition of done for every increment was the same and is worth stating because
it explains the shape of the repository: \emph{the feature works end to end, its
service-layer lines are covered by tests, the relevant document is updated, and
\code{mvn verify} is green}. Documentation was treated as part of the increment, not
as a final phase, which is why the API reference, BLE protocol and firmware notes
match the code.

Waterfall was rejected because the hardware findings would have arrived in the
verification phase, after the backend design had been frozen. A pure prototype
model was rejected because the deliverable is an assessed project with a written
specification, a documented threat model and reproducible evidence; those artefacts
need a plan that produces them incrementally.

\section{Project Schedule}

Development ran for eleven weeks, with a short baseline week at the start and a
final verification week. Sprints were sequenced by dependency: identity first, because
every other endpoint needs a caller; money second, because the exit gate depends on
payment; hardware third, once the contract was stable; verification and writing
last, against a system that no longer changed shape.

\begin{table}[htbp]
  \centering
  \small
  \begin{tabularx}{\textwidth}{@{}llX@{}}
    \toprule
    \textbf{Sprint} & \textbf{Dates} & \textbf{Goal and deliverable} \\
    \midrule
    S0 & 07--11 Aug & Baseline and decisions. Repository, architecture
      specification, decision log D1--D23, risk list. \\
    S1 & 12--18 Aug & Identity. Registration, login, JWT issue and verification,
      security filter chain. \\
    S2 & 19--25 Aug & Vehicles and facilities. Owner-scoped CRUD, facility lookup by
      code, admin-key surface, demo seeder. \\
    S3 & 26 Aug--01 Sep & Sessions and money. Creation, anti-abuse rule, per-second
      fee, mock payment, cancellation, state machine. \\
    S4 & 02--08 Sep & Gate authorisation and devices. Entry/exit verdicts, device
      provisioning and secrets, remote-command queue. \\
    S5 & 09--15 Sep & Lifecycle hardening. Stale-session sweep, grace window,
      exception advice, history paging. \\
    S6 & 16--22 Sep & Mobile application. Typed REST client, token persistence,
      eleven screens on real data. \\
    S7 & 23--29 Sep & BLE. Discovery, GATT round trip, real transport in the
      application; firmware transports and servo. \\
    S8 & 30 Sep--08 Oct & Firmware completion. Remote-command poller, stall
      watchdog, serial bridge for the demo. \\
    S9 & 09--15 Oct & Verification. Docker Compose test workflow, coverage gate,
      report collection. \\
    S10 & 16--22 Oct & Documentation and hand-over. API reference, protocol and
      hardware guides, Postman collection, this report. \\
    \bottomrule
  \end{tabularx}
  \caption{Sprint plan.}
  \label{tab:sprints}
\end{table}

\reportimage[0.9]{\ImgGantt}{Gantt chart of the ten-week schedule, showing the
dependency between identity, billing, hardware and verification work.}{fig:gantt}

Sprints S1--S5 built the backend and its tests; S6--S8 built the phone and the gate
against that contract; S9 built the machinery that proves the whole thing works. The
ordering meant that the firmware was written against an API that had already been
tested, so a gate failure could be attributed to hardware or transport rather than
to an unfinished backend.

\section{Risk Analysis}

Risks were reviewed at the start of each sprint. Table~\ref{tab:risks} lists the
significant ones, the mitigation chosen and what actually happened.

\begin{table}[htbp]
  \centering
  \small
  \begin{tabularx}{\textwidth}{@{}XllX@{}}
    \toprule
    \textbf{Risk} & \textbf{Likelihood} & \textbf{Impact} &
    \textbf{Mitigation and outcome} \\
    \midrule
    Servo browns out the ESP32 when the gate moves & High & High &
      External 5\,V supply, common ground and a decoupling capacitor were made
      mandatory; the actuator is behind an interface so an LED stand-in can be
      swapped in one line. The servo build is current. \\
    BLE notification lost because the phone subscribed after the write & Medium &
      High & The client subscribes to the status characteristic before writing and
      filters notifications by session id. Implemented. \\
    Backend and app disagree about the response shape & High & Medium &
      One envelope for every application endpoint, two documented exceptions for
      firmware, and a typed client that rejects unknown shapes. \\
    Double payment through a double tap or retry & Medium & High &
      Session row locked during payment; the second transaction sees
      \code{PAYMENT\_PENDING} and is refused. \\
    Concurrent session creation creates two sessions & Low & High &
      The user row is locked for the transaction, serialising the read-then-write
      check. \\
    Coverage gate blocks iteration & Medium & Medium &
      Tests are written with the feature, not afterwards; the gate is scoped to the
      service layer so generated code does not dilute it. \\
    Hardware unavailable for testing & High & Medium &
      The serial-bridge transport allows the full backend flow without Wi-Fi; BLE
      and servo behaviour remain documented as bench steps. \\
    Secrets committed to the repository & Medium & High &
      \code{.env} is ignored, sample values live in \code{.env.example}, device
      secrets are hashed, and the admin key has no default. \\
    Scope creep from payment or occupancy features & Medium & Medium &
      Both were written into the non-goals and revisited only as future work. \\
    Test environment differs from deployment & Low & High &
      No H2 and no Testcontainers: tests run against PostgreSQL 16, the same image
      the stack uses. \\
    \bottomrule
  \end{tabularx}
  \caption{Risk register with the mitigation that was adopted.}
  \label{tab:risks}
\end{table}

\section{Effort Estimation}

A simple decomposition estimate was used: each module was sized by the number of
distinct behaviours it had to support, and reviewed at the end of the sprint in
which it was built. Table~\ref{tab:effort} records the planned effort and the
measured size of the delivered code. The size figures are counted from the
repository; the effort figures are planning estimates for a three-person team
working part-time.

\begin{table}[htbp]
  \centering
  \small
  \begin{tabularx}{\textwidth}{@{}lrrX@{}}
    \toprule
    \textbf{Work package} & \textbf{Planned} & \textbf{Size} & \textbf{Notes} \\
    & \textbf{person-days} & \textbf{(lines)} & \\
    \midrule
    Requirements and design & 6 & --- & specification, decision log, threat model \\
    Backend domain and services & 14 & 3\,275 & Java, entities to controllers \\
    Backend tests & 8 & 2\,240 & 15 classes, 118 tests, Mockito \\
    Mobile application & 14 & 2\,278 & TypeScript, 11 screens, REST and BLE \\
    Gate firmware & 10 & 1\,761 & Arduino, transports, actuator, poller \\
    Build and test workflow & 6 & 261 & Compose profile and two shell scripts \\
    Documentation and report & 8 & --- & six guides, API reference, this report \\
    \midrule
    \textbf{Total} & \textbf{66} & \textbf{9\,815} & \\
    \bottomrule
  \end{tabularx}
  \caption{Indicative effort distribution. Size counts exclude blank lines and
  generated build output; scripts are shell, firmware is C/C++.}
  \label{tab:effort}
\end{table}

The distribution is deliberately backend-heavy. The mobile application and the
firmware are thin by design: neither holds a decision that the backend could make.
The largest single risk-reduction item was not a feature but the test workflow,
because it turned every later change into a verifiable one.
